% Options for packages loaded elsewhere
\PassOptionsToPackage{unicode}{hyperref}
\PassOptionsToPackage{hyphens}{url}
\PassOptionsToPackage{dvipsnames,svgnames,x11names}{xcolor}
%
\documentclass[
  10pt,
  a4paper,
]{article}
\usepackage{amsmath,amssymb}
\usepackage{iftex}
\ifPDFTeX
  \usepackage[T1]{fontenc}
  \usepackage[utf8]{inputenc}
  \usepackage{textcomp} % provide euro and other symbols
\else % if luatex or xetex
  \usepackage{unicode-math} % this also loads fontspec
  \defaultfontfeatures{Scale=MatchLowercase}
  \defaultfontfeatures[\rmfamily]{Ligatures=TeX,Scale=1}
\fi
\usepackage{lmodern}
\ifPDFTeX\else
  % xetex/luatex font selection
\fi
% Use upquote if available, for straight quotes in verbatim environments
\IfFileExists{upquote.sty}{\usepackage{upquote}}{}
\IfFileExists{microtype.sty}{% use microtype if available
  \usepackage[]{microtype}
  \UseMicrotypeSet[protrusion]{basicmath} % disable protrusion for tt fonts
}{}
\makeatletter
\@ifundefined{KOMAClassName}{% if non-KOMA class
  \IfFileExists{parskip.sty}{%
    \usepackage{parskip}
  }{% else
    \setlength{\parindent}{0pt}
    \setlength{\parskip}{6pt plus 2pt minus 1pt}}
}{% if KOMA class
  \KOMAoptions{parskip=half}}
\makeatother
\usepackage{xcolor}
\usepackage[margin=23mm]{geometry}
\usepackage{color}
\usepackage{fancyvrb}
\newcommand{\VerbBar}{|}
\newcommand{\VERB}{\Verb[commandchars=\\\{\}]}
\DefineVerbatimEnvironment{Highlighting}{Verbatim}{commandchars=\\\{\}}
% Add ',fontsize=\small' for more characters per line
\newenvironment{Shaded}{}{}
\newcommand{\AlertTok}[1]{\textcolor[rgb]{1.00,0.00,0.00}{\textbf{#1}}}
\newcommand{\AnnotationTok}[1]{\textcolor[rgb]{0.38,0.63,0.69}{\textbf{\textit{#1}}}}
\newcommand{\AttributeTok}[1]{\textcolor[rgb]{0.49,0.56,0.16}{#1}}
\newcommand{\BaseNTok}[1]{\textcolor[rgb]{0.25,0.63,0.44}{#1}}
\newcommand{\BuiltInTok}[1]{\textcolor[rgb]{0.00,0.50,0.00}{#1}}
\newcommand{\CharTok}[1]{\textcolor[rgb]{0.25,0.44,0.63}{#1}}
\newcommand{\CommentTok}[1]{\textcolor[rgb]{0.38,0.63,0.69}{\textit{#1}}}
\newcommand{\CommentVarTok}[1]{\textcolor[rgb]{0.38,0.63,0.69}{\textbf{\textit{#1}}}}
\newcommand{\ConstantTok}[1]{\textcolor[rgb]{0.53,0.00,0.00}{#1}}
\newcommand{\ControlFlowTok}[1]{\textcolor[rgb]{0.00,0.44,0.13}{\textbf{#1}}}
\newcommand{\DataTypeTok}[1]{\textcolor[rgb]{0.56,0.13,0.00}{#1}}
\newcommand{\DecValTok}[1]{\textcolor[rgb]{0.25,0.63,0.44}{#1}}
\newcommand{\DocumentationTok}[1]{\textcolor[rgb]{0.73,0.13,0.13}{\textit{#1}}}
\newcommand{\ErrorTok}[1]{\textcolor[rgb]{1.00,0.00,0.00}{\textbf{#1}}}
\newcommand{\ExtensionTok}[1]{#1}
\newcommand{\FloatTok}[1]{\textcolor[rgb]{0.25,0.63,0.44}{#1}}
\newcommand{\FunctionTok}[1]{\textcolor[rgb]{0.02,0.16,0.49}{#1}}
\newcommand{\ImportTok}[1]{\textcolor[rgb]{0.00,0.50,0.00}{\textbf{#1}}}
\newcommand{\InformationTok}[1]{\textcolor[rgb]{0.38,0.63,0.69}{\textbf{\textit{#1}}}}
\newcommand{\KeywordTok}[1]{\textcolor[rgb]{0.00,0.44,0.13}{\textbf{#1}}}
\newcommand{\NormalTok}[1]{#1}
\newcommand{\OperatorTok}[1]{\textcolor[rgb]{0.40,0.40,0.40}{#1}}
\newcommand{\OtherTok}[1]{\textcolor[rgb]{0.00,0.44,0.13}{#1}}
\newcommand{\PreprocessorTok}[1]{\textcolor[rgb]{0.74,0.48,0.00}{#1}}
\newcommand{\RegionMarkerTok}[1]{#1}
\newcommand{\SpecialCharTok}[1]{\textcolor[rgb]{0.25,0.44,0.63}{#1}}
\newcommand{\SpecialStringTok}[1]{\textcolor[rgb]{0.73,0.40,0.53}{#1}}
\newcommand{\StringTok}[1]{\textcolor[rgb]{0.25,0.44,0.63}{#1}}
\newcommand{\VariableTok}[1]{\textcolor[rgb]{0.10,0.09,0.49}{#1}}
\newcommand{\VerbatimStringTok}[1]{\textcolor[rgb]{0.25,0.44,0.63}{#1}}
\newcommand{\WarningTok}[1]{\textcolor[rgb]{0.38,0.63,0.69}{\textbf{\textit{#1}}}}
\setlength{\emergencystretch}{3em} % prevent overfull lines
\providecommand{\tightlist}{%
  \setlength{\itemsep}{0pt}\setlength{\parskip}{0pt}}
\setcounter{secnumdepth}{-\maxdimen} % remove section numbering
\ifLuaTeX
  \usepackage{selnolig}  % disable illegal ligatures
\fi
\usepackage{bookmark}
\IfFileExists{xurl.sty}{\usepackage{xurl}}{} % add URL line breaks if available
\urlstyle{same}
\hypersetup{
  colorlinks=true,
  linkcolor={black},
  filecolor={Maroon},
  citecolor={Blue},
  urlcolor={blue},
  pdfcreator={LaTeX via pandoc}}

\author{}
\date{}

\begin{document}

\section{DDTA research work plan after documentation-layer
closure}\label{ddta-research-work-plan-after-documentation-layer-closure}

\textbf{WORK PLAN - REVISION 20}

\textbf{Status:} active BA5 execution plan after BA5-T1 canonical
referent naming and controlled-authoring pressure testing.

\textbf{Supersedes:} Revision 19 only for forward execution state;
R1-R19 remain historical research records.

\subsection{1. Fixed prior state}\label{fixed-prior-state}

\begin{itemize}
\tightlist
\item
  Chapters 2-4: CLOSED / FINAL for current thesis scope.
\item
  Documentation layer: CLOSED.
\item
  BA0 responsibility/non-goals: CLOSED.
\item
  BA1 minimal BAE identity ontology: CLOSED.
\item
  BA2 relation/action vocabulary: CLOSED BY BA2-T4.
\item
  BA3 provenance/derivation/identity/lifecycle/change-revalidation:
  CLOSED BY BA3-T4.
\item
  BA4 projection boundary/traceability/coverage/interpretation: CLOSED
  BY BA4-T3.
\item
  \texttt{BAReferent} and \texttt{BAProposition}: ACCEPTED first-class
  semantic identity families.
\item
  ThreatForge remains a case study/tooling instrument, not DDTA semantic
  authority.
\end{itemize}

\subsection{2. BA5 objective - canonical semantic registry and
controlled
authoring}\label{ba5-objective---canonical-semantic-registry-and-controlled-authoring}

BA5 now starts from the stronger hypothesis that semantic authoring
should be canonical by construction rather than synonym-normalized after
the fact.

For named project referents and, subject to later pressure, the wider
operative BA vocabulary:

\begin{Shaded}
\begin{Highlighting}[]
\NormalTok{registered canonical token}
\NormalTok{    {-}\textgreater{} normative semantic authoring}

\NormalTok{unregistered alias / synonym}
\NormalTok{    {-}\textgreater{} not silently equivalent}
\end{Highlighting}
\end{Shaded}

The purpose is to maximize cross-document and cross-view consistency
before investing in lexical normalization machinery.

NLP/LLM assistance remains downstream candidate-producing help and is
not required by the current core hypothesis.

\subsection{3. BA5-T1 result}\label{ba5-t1-result}

\textbf{Status: COMPLETED / PROVISIONAL PASS WITH
CANONICAL-REFERENT-NAMING LOWER-BOUND.}

T1 establishes provisionally:

\begin{Shaded}
\begin{Highlighting}[]
\NormalTok{one named referent {-}\textgreater{} one exact canonicalName per baseline/scope   REQUIRED}
\NormalTok{same referent same name across governed semantic references        REQUIRED}
\NormalTok{same referent same name across shared derived views                REQUIRED}
\NormalTok{alias/synonym as second normative entity identifier                REJECTED}
\NormalTok{case/format variant as implicit alias                              REJECTED}
\NormalTok{canonical{-}name collision in one naming scope                       REJECTED}
\NormalTok{canonicalName == BAReferent identity                                REJECTED}
\NormalTok{governed cross{-}baseline rename + BA3 RETAIN                         PASS}
\NormalTok{projection{-}owned type/category labels                              ALLOWED}
\NormalTok{method{-}owned non{-}referent labels                                   ALLOWED DOWNSTREAM}
\NormalTok{tool exact validation/completion                                   ALLOWED}
\NormalTok{tool hidden alias normalization                                    REJECTED IN CORE}
\end{Highlighting}
\end{Shaded}

The decisive refinement is that naming invariance is baseline-scoped. A
governed rename does not itself force new \texttt{BAReferent} identity.

\subsection{4. Why this does not reopen
BA1-BA4}\label{why-this-does-not-reopen-ba1-ba4}

\begin{itemize}
\tightlist
\item
  BA1 identity remains semantic; canonical name is lexical metadata.
\item
  BA2 already distinguishes stable semantic keys from lexical wording
  and needs no \texttt{nameOf} operator.
\item
  BA3 already supports RETAIN under wording/name change when meaning
  survives.
\item
  BA4 projection-local taxonomy remains free; only the shared referent
  name is constrained when a projection presents that referent.
\end{itemize}

No new BAE family or BA2 operator is forced.

\subsection{5. Registry responsibility}\label{registry-responsibility}

BA5-T1 requires an auditable governed canonical-name registry
responsibility for named referents, but does not mandate one physical
implementation.

Minimum properties:

\begin{itemize}
\tightlist
\item
  baseline/scope binding;
\item
  exact canonical token;
\item
  unambiguous referent resolution;
\item
  uniqueness within the active naming scope;
\item
  no normative alias equivalence by default.
\end{itemize}

A tool may enforce the registry, but cannot become semantic authority by
inventing equivalence or accepted terms.

\subsection{6. BA5-T2 - canonical semantic registry coverage and
governed-extension pressure
test}\label{ba5-t2---canonical-semantic-registry-coverage-and-governed-extension-pressure-test}

\textbf{NEXT / NOT EXECUTED BY THIS PLAN REVISION.}

BA5-T2 must pressure-test whether the same controlled-registry
discipline can cover the wider operative BA vocabulary already
established by BA2 without introducing unnecessary lexical machinery.

It must test at least:

\begin{enumerate}
\def\labelenumi{\arabic{enumi}.}
\tightlist
\item
  the thirteen accepted \texttt{semanticOperatorKey} values;
\item
  operator-scoped \texttt{roleKey} values;
\item
  method-neutral \texttt{semanticKind} keys, including current evidence
  such as \texttt{store} and \texttt{contract};
\item
  controlled values such as responsibility kinds and other
  non-identifiable semantic values;
\item
  exact-token validation and registry lookup across documentation
  -\textgreater{} BA -\textgreater{} projections;
\item
  the minimum governed workflow for admitting a genuinely new canonical
  term;
\item
  attempted alias/synonym introduction as a negative control; and
\item
  whether tool completion/validation can operate without
  semantic-authority leakage.
\end{enumerate}

\subsubsection{BA5-T2 guardrails}\label{ba5-t2-guardrails}

Do not:

\begin{itemize}
\tightlist
\item
  build a synonym ontology by default;
\item
  perform arbitrary-narrative migration;
\item
  authorize NLP/LLM output as semantic truth;
\item
  convert method-specific projection taxonomy into BA vocabulary;
\item
  create new BAE families merely to store lexical keys;
\item
  reopen BA0-BA4 without a material semantic counterexample;
\item
  start BA6.
\end{itemize}

\subsubsection{BA5-T2 exit condition}\label{ba5-t2-exit-condition}

Produce the next smallest controlled-registry lower-bound refinement and
identify whether any concrete evidence actually forces synonym/alias
support.

Do not close BA5 unless the integrated lexical/registry responsibilities
have survived sufficient pressure.

\subsection{7. BA6 - complete Base Analysis regression and
closure}\label{ba6---complete-base-analysis-regression-and-closure}

Only after BA5 closes, regress the complete BA0-BA5 design against the
closed corpora and at least one structurally different holdout.

BA6 should emphasize the complete integration path:

\begin{Shaded}
\begin{Highlighting}[]
\NormalTok{governed documentation}
\NormalTok{ {-}\textgreater{} canonical registry{-}controlled authoring}
\NormalTok{ {-}\textgreater{} accepted BA}
\NormalTok{ {-}\textgreater{} provenance/change continuity}
\NormalTok{ {-}\textgreater{} multiple projections}
\NormalTok{ {-}\textgreater{} rebuild / cross{-}projection comparison}
\end{Highlighting}
\end{Shaded}

BA6 may reopen only the smallest earlier responsibility on a material
counterexample.

Only BA6 may close the complete Base Analysis milestone for current
thesis scope.

\subsection{8. Analysis envelope remains
downstream}\label{analysis-envelope-remains-downstream}

Only after Base Analysis closure:

\begin{itemize}
\tightlist
\item
  A1 - AnalysisRecord / execution envelope;
\item
  A2 - common reviewed Finding boundary;
\item
  A3 - change/provenance integration;
\item
  formal methodology overlays and STRIDE/STRIDE-AI evaluations.
\end{itemize}

\subsection{9. Next authorized
microstep}\label{next-authorized-microstep}

Only after the BA5-T1 package is reviewed, committed, pushed and
remotely verified, execute:

\begin{quote}
\textbf{BA5-T2 - canonical semantic registry coverage and
governed-extension pressure test.}
\end{quote}

Do not start BA6 or downstream analysis schemas before BA5 is explicitly
closed.

\end{document}
